\section{The geometry of unresolved prior information}
\label{sec:directional-geometry}
The acquired-history image and uncertainty about the prior are different
geometric objects. The first has its dimension bounded by $n(r-1)$.
The second can occupy every future direction, even after a fixed history.
We describe the second object here. This also identifies the directions
which are suppressed when an uncertainty problem is summarized by a
single squared-error floor.

\subsection{The bounded-test principle}
The finite-radius mechanism in this section is independent of monomials.
We state it before its confluent specialization. This formulation and the
normalization below were isolated in the preceding mathematical audit
\cite{A1v14audit}; we retain that attribution. The result is a bounded
right-inverse principle, not a second positive-history attainment theorem.

\begin{proposition}[Relative prior perturbations under an observation map]
\label{prop:general-bounded-dual}
Let $\mu$ be a probability, $0<c\le\ell\le C<\infty$, and
$\nu=\ell\mu/(\mu\ell)$. Let $\phi=(\phi_1,\ldots,\phi_q)$ be
bounded real tests with $\|\phi_j\|_\infty\le B$ and
$\operatorname{Cov}_\nu(\phi,\phi)\ge\lambda I_q$, $\lambda>0$.
For $0<\varepsilon\le1/2$, vary $\mu'$ over
\[
 \{(1+u)\mu:\mu u=0,\ \|u\|_\infty\le\varepsilon\},
 \qquad \nu'=\ell\mu'/(\mu'\ell).
\]
For every linear map $A:\R^q\to\R^d$ the displacement set satisfies
\[
 c_0\varepsilon A[-1,1]^q
 \ \subset\ \{A(\nu'\phi-\nu\phi)\}
 \ \subset\ C_0\varepsilon A[-1,1]^q,
\]
where $c_0,C_0>0$ depend only on $q,B,\lambda$, not on the singular
values or rank of $A$.
\end{proposition}
\begin{proof}
For $\mu'=(1+u)\mu$, the exact identity
\[
 \nu'\phi-\nu\phi
 =\frac{\nu\{u(\phi-\nu\phi)\}}{1+\nu u}
\]
bounds each coordinate by $4B\varepsilon$. For the other inclusion put
$\Gamma=\operatorname{Cov}_\nu(\phi,\phi)$ and, for
$v\in[-1,1]^q$, set
\[
 f_v=v^{\mathsf T}\Gamma^{-1}(\phi-\nu\phi),\qquad
 K=\max\{1,2qB/\lambda\},\qquad s=\varepsilon/(4K).
\]
Then $\nu f_v=0$, $\nu(\phi f_v)=v$, and $\|f_v\|_\infty\le K$.
Define a prior, rather than only a posterior perturbation, by
\[
 \frac{d\mu'}{d\mu}=\frac{1+s f_v}{1+s\mu f_v}.
\]
It has total mass one, is positive, and obeys
\[
 \left\|\frac{d\mu'}{d\mu}-1\right\|_\infty
 \le\frac{2sK}{1-sK}
 =\frac{\varepsilon/2}{1-\varepsilon/4}\le\varepsilon.
\]
The evidence is $\mu\ell/(1+s\mu f_v)$, so Bayes normalization gives
exactly $\nu'=(1+s f_v)\nu$ and hence displacement $sv$. This realizes
one common cube, not merely separate directional alternatives. Applying
$A$ to both inclusions completes the proof, also when $A$ has a kernel.
\end{proof}

\subsection{Complete confluent flags and all width orders}


Throughout this section the calibration is known. Write $q=q_m$ and use
one of the Leja orders in Lemma~\ref{lem:leja-scales}. Put
\[
 \phi_a=(g_1^a,\ldots,g_q^a),\qquad
 D_a=\operatorname{diag}(d_1,\ldots,d_q),\qquad
 B_\infty^q=[-1,1]^q,
\]
and let $L_a$ be the matrix in \eqref{eq:leja-physical}. All repeated
formal labels are retained. The convention $d_{q+1}=0$ will be useful.
For a probability $\nu$, let
\[
 X_a(\nu)=(\nu(t^{Hx_i}))_{i=1}^q=L_aD_a\nu\phi_a.
\]
The norm of a difference of physical query vectors is equivalent to the
Euclidean norm of the corresponding difference of $X_a$, with the fixed
constants already established for $G_m$.

Fix a full-support prior $\mu$ and a history $h$ of length $n$, where
$n+m\le N$. Its likelihood is denoted by $\ell_h$, and its posterior
by $\nu_h=\ell_h\mu/\mu\ell_h$. For $0<\varepsilon\le1/2$ set
\begin{equation}\label{eq:relative-prior-ball}
 \mathcal B_\varepsilon(\mu)=
 \{(1+u)\mu:\ \mu u=0,\ \|u\|_{L^\infty(\mu)}\le\varepsilon\}.
\end{equation}
Every member is a full-support probability. Let $\nu_h'$ denote the
posterior obtained from $\mu'\in\mathcal B_\varepsilon(\mu)$ using
\emph{the same} likelihood, and define the centered ambiguity set
\begin{equation}\label{eq:ambiguity-body}
 \mathcal A_h(a,\varepsilon)=
 \{X_a(\nu_h')-X_a(\nu_h):\mu'\in\mathcal B_\varepsilon(\mu)\}.
\end{equation}
This is a set at a specified history; choosing a point in it does not
select a different prior after each realized history in a statistical
experiment. A separate common-prior decision statement is given below.

\begin{lemma}[A uniformly bounded dual for the complete test flag]
\label{lem:flag-covariance}
The covariance matrix
\[
 \Gamma_{a,h}=\operatorname{Cov}_{\nu_h}(\phi_a,\phi_a)
\]
has least eigenvalue bounded below by a positive constant, uniformly in
$a\in K$, histories of length at most $N-m$, and all orders of the
formal nodes, including their repeated-node limits. Moreover there is
$C_0\ge1$ such that, for every $v\in B_\infty^q$, the function
\begin{equation}\label{eq:bounded-flag-dual}
 f_v(t)=v^{\mathsf T}\Gamma_{a,h}^{-1}
                 (\phi_a(t)-\nu_h\phi_a)
\end{equation}
satisfies $\nu_h f_v=0$, $\nu_h(\phi_a f_v)=v$, and
$\|f_v\|_\infty\le C_0$. The constants can be chosen uniformly for
centers in any fixed dominated prior family
$c_-\mu_0\le\mu\le c_+\mu_0$ with $c_->0$ and full-support $\mu_0$.
\end{lemma}
\begin{proof}
For every ordered multiset of $q$ positive nodes, the functions
$1,g_1^a,\ldots,g_q^a$ are linearly independent. Indeed,
Lemma~\ref{lem:prefix-hermite} identifies their span with the constant
and complete positive-exponent Hermite blocks. Linear independence of
these exponential-polynomial functions on $(0,1)$ follows from the
confluent Chebyshev argument in Lemma~\ref{lem:confluent-positive}.
A nonzero linear combination of the $g_j^a$ cannot therefore be constant
on the support of a full-support measure. Consequently its variance is
positive, and the covariance matrix under $\mu_0$ is positive definite.

The uniform bounds and continuity in the proof of
Lemma~\ref{lem:newton-attainment} apply to the full prefix of length
$q$, regardless of the number of distinct node values. They show that
this covariance matrix is continuous on $K$ for each fixed permutation.
Compactness, and the finite number of permutations, give
$\operatorname{Cov}_{\mu_0}(\phi_a,\phi_a)\ge\lambda_0 I_q$
for some $\lambda_0>0$. This statement concerns the normalized functions,
not physically observable coordinates divided by a zero pivot.

Since $\kappa^n\le\ell_h\le1$, one has
$\nu_h\ge\kappa^n c_-\mu_0$. For any $w\in\R^q$, the variational
formula for variance gives
\[
 \inf_b\nu_h((w^{\mathsf T}\phi_a-b)^2)
 \ge\kappa^n c_-\inf_b\mu_0((w^{\mathsf T}\phi_a-b)^2)
 \ge\kappa^N c_-\lambda_0\|w\|^2.
\]
For a single fixed prior use $\mu_0=\mu$ and $c_-=c_+=1$.
This proves the eigenvalue assertion without any compactness assumption
on the family of priors. Formula~\eqref{eq:bounded-flag-dual} now has
zero mean, and multiplication by the centered test vector gives $v$.
The uniform bound on $\phi_a$ and the just obtained bound on
$\Gamma_{a,h}^{-1}$ give $C_0$, after increasing it to at least one.
\end{proof}

\begin{proposition}[All collision scales of the prior ambiguity body]
\label{prop:ambiguity-ellipsoid}
There are constants $c,C>0$, with the uniformity of
Lemma~\ref{lem:flag-covariance}, such that
\begin{equation}\label{eq:ambiguity-sandwich}
 c\varepsilon L_aD_aB_\infty^q
 \ \subset\ \mathcal A_h(a,\varepsilon)
 \ \subset\ C\varepsilon L_aD_aB_\infty^q.
\end{equation}
If
\[
 w_k(\mathcal A)=\inf_{\dim V\le k}\sup_{z\in\mathcal A}
                        \operatorname{dist}(z,V),\qquad 0\le k\le q,
\]
then
\begin{equation}\label{eq:ambiguity-widths}
 c\varepsilon d_{k+1}\le w_k(\mathcal A_h(a,\varepsilon))
                         \le C\varepsilon d_{k+1}.
\end{equation}
In particular all widths past the number of distinct positive future
exponents are zero. No past-dimension truncation occurs in this assertion.
\end{proposition}
\begin{proof}
Write $\nu=\nu_h$. For $\mu'=(1+u)\mu$ in
\eqref{eq:relative-prior-ball}, direct normalization yields
\[
 \nu'\phi_a-\nu\phi_a
 =\frac{\nu\bigl(u(\phi_a-\nu\phi_a)\bigr)}{1+\nu u}.
\]
If every $g_j^a$ is bounded by $B$, the right side has coordinatewise
absolute value at most $2B\varepsilon/(1-\varepsilon)\le4B\varepsilon$.
Multiplication by $L_aD_a$ proves the upper inclusion.

For the lower inclusion take any $v\in B_\infty^q$, let $f=f_v$ be
\eqref{eq:bounded-flag-dual}, and put $s=\varepsilon/(4C_0)$. Define
\begin{equation}\label{eq:pullback-posterior-tilt}
 d\mu' =\frac{1+s f}{1+s\mu f}\,d\mu.
\end{equation}
This is a positive probability, and
\[
 \left\|\frac{d\mu'}{d\mu}-1\right\|_\infty
 \le\frac{2s C_0}{1-s C_0}
 =\frac{\varepsilon/2}{1-\varepsilon/4}\le\varepsilon.
\]
Because $\nu f=0$, its posterior at the same history is exactly
$d\nu'=(1+s f)\,d\nu$. Hence
$\nu'\phi_a-\nu\phi_a=s v$, proving the lower inclusion with
$c=1/(4C_0)$. Formula~\eqref{eq:pullback-posterior-tilt} enforces prior
normalization; simply using $(1+s f)\mu$ here would not do so.

Let $s_a$ be the number of distinct positive future exponents. The active
$q$ by $s_a$ block of $L_a$ is bounded above and below as an operator,
by its leading triangular block in Lemma~\ref{lem:leja-scales}.
Thus the nonzero singular values of
$L_a\operatorname{diag}(d_1,\ldots,d_{s_a})$ are comparable, in order,
to $d_1,\ldots,d_{s_a}$. The comparison follows directly by applying
the min--max formula to the two quadratic forms, bounded above and below
by fixed multiples of $\sum_jd_j^2v_j^2$.
The Euclidean unit ball lies in the cube, which lies in the ball of
radius $\sqrt q$. An ellipsoid with ordered semiaxes $b_j$ has $k$th
linear width $b_{k+1}$: projection on its first $k$ principal axes gives
the upper bound, while any $k$-dimensional space has a nonzero
orthogonal direction in the span of its first $k+1$ axes, giving the
lower bound. Apply this fact to the two inclusions in
\eqref{eq:ambiguity-sandwich}. If $k\ge s_a$, all vectors lie in a
space of dimension $s_a$, and both sides of
\eqref{eq:ambiguity-widths} are zero.
\end{proof}

\subsection{Exact lower moments and a vanishing uncertainty direction}
The preceding proposition resolves every direction, rather than only
the longest one. An equivalent formulation specifies which moments have
already been supplied. Let $s_a$ again denote the number of distinct
positive nodes. For $0\le k\le q$, restrict
\eqref{eq:relative-prior-ball} by the finite conditions
\begin{equation}\label{eq:prefix-advice}
 \nu_h'(t^{Hx_i})=\nu_h(t^{Hx_i}),
                 \qquad 1\le i\le\min\{k,s_a\},
\end{equation}
and call the restricted class $\mathcal B_{h,k}$. Denote by
$Q_m^a(\nu)$ the vector of probabilities of the ordered physical
$m$-trial queries, with its normalized Euclidean norm $\|\cdot\|_{\rm av}$.
Define its squared minimax radius by
\[
 \mathfrak a_{h,k}(a,\varepsilon)=
 \inf_z\sup_{\mu'\in\mathcal B_{h,k}}
       \|z-Q_m^a(\nu_h')\|_{\rm av}^2.
\]
The center $\mu$, the relative envelope, the calibration and the values
in \eqref{eq:prefix-advice} specify this local problem. These are exact
conditional-moment constraints at $h$, not the pre-acquisition,
absolute-error moment name of Section~\ref{sec:uncertainty-geometry}.

\begin{corollary}[Sharp residual uncertainty after a moment prefix]
\label{cor:prefix-ambiguity}
With uniform constants as above,
\begin{equation}\label{eq:prefix-ambiguity-risk}
 c\varepsilon^2d_{k+1}^2
 \le\mathfrak a_{h,k}(a,\varepsilon)
 \le C\varepsilon^2d_{k+1}^2.
\end{equation}
For a constant-failure history, the same order is realized as an
unconditional, event-weighted physical prediction problem with exact
prior-moment advice. In particular, for
$A_\theta=\{0,1,2+\theta\}$, $N=3$, $n=1$, $m=2$ and sufficiently
small $|\theta|$, one has
\begin{equation}\label{eq:invisible-memory-visible-ambiguity}
 \Xi_3(M,A_\theta)\asymp M^{-1},\qquad
 \mathfrak a_{h,4}(A_\theta,\varepsilon)
                            \asymp\varepsilon^2\theta^2.
\end{equation}
The latter is zero at $\theta=0$.
\end{corollary}
\begin{proof}
For $k<s_a$, the first $k$ raw constraints in
\eqref{eq:prefix-advice} are equivalent to equality of the first $k$
normalized moments. This follows by triangular back substitution in
the leading block of $L_aD_a$, whose first $k$ pivots are nonzero.
This equivalence is an exact identity, not a stability assertion for
approximate constraints. The upper-inclusion proof in
Proposition~\ref{prop:ambiguity-ellipsoid} therefore has its first $k$
coordinates zero. Its remaining scaled coordinates are bounded by
$C\varepsilon d_{k+1}$. The central prediction gives the upper bound.
For the lower bound use the construction
\eqref{eq:pullback-posterior-tilt} with $v=e_{k+1}$ and $v=-e_{k+1}$.
Both priors satisfy the constraints and their query vectors differ by
at least $c\varepsilon d_{k+1}$. For any proposed output, one of its
two distances is at least half this separation. Query-norm equivalence
proves the lower bound. If $k\ge s_a$, all distinct raw future moments
are fixed, and Proposition~\ref{prop:tests} gives radius zero.

For the physical assertion issue $n$ constant commands whose failure
factor is $1/2$, and consider the event $E$ of $n$ failures. Its
likelihood is the constant $2^{-n}$, and its probability is $2^{-n}$
under every prior in the class. On $E$ the posterior is the prior, so
\eqref{eq:prefix-advice} is now exact \emph{prior}-moment advice fixed
before acquisition. Score the squared prediction loss for the selected
$m$-trial query on $E$, and assign zero loss off $E$. Every trial and
the probability of $E$ remain counted. A single rule has the same
prediction kernel on $E$ for all compatible priors. The preceding
two-point argument, multiplied by $2^{-n}$, gives the unconditional
lower bound; the central prediction gives the upper bound. This is a
specific event-weighted decision problem, not a claim about the risk
on all other report histories under a different payoff.

For the displayed family, the five nonconstant formal two-trial nodes,
before the fixed normalization by $H$, are
\[
 1,\quad 2,\quad 2+\theta,\quad 3+\theta,\quad4+2\theta.
\]
They form four uniformly separated groups, with the only merging pair
$2,2+\theta$. Thus $\mathcal V_{2,j}\asymp1$ for $1\le j\le4$ and
$\mathcal V_{2,5}\asymp|\theta|$. Lemma~\ref{lem:leja-scales}
gives $d_1,\ldots,d_4\asymp1$ and $d_5\asymp|\theta|$,
including $d_5=0$ at the collision. Formula~\eqref{eq:prefix-ambiguity-risk}
proves the second assertion. At checkpoint $(1,2)$ the attainable
truncation is $p_{1,2}=2$, so its regret profile is $\asymp M^{-1}$.
At checkpoint $(2,1)$ there are two uniformly separated positive
one-trial nodes, and the same profile results. Taking the maximum proves
the first assertion in \eqref{eq:invisible-memory-visible-ambiguity}.
\end{proof}

The example separates a collision which is invisible to the acquired
memory exponent at this horizon from a vanishing direction of unresolved
prior information. More generally, \eqref{eq:ambiguity-widths} uses
all future scales, whereas \eqref{eq:intrinsic-profile} uses only those
which the past can excite. The isotropic common-name law in the next
section allows uncertainty in a uniformly visible direction; its
$\delta^2$ term is therefore not, by itself, a classification of the
weaker directions. No calibration-only converse is inferred from the
prior perturbations constructed here.
